By Lemma~\(\mathrm{lem{:}clw\text{-}population\text{-}design\text{-}rule}\), the balanced-treatment population criterion of Carneiro--Lee--Wilhelm is linear-regression residual variance divided by the budget-feasible sample size.  Definition~\(\mathrm{def{:}clw\text{-}linear\text{-}cost\text{-}selector}\) transports only that criterion and specializes the cost to \(nC_d\), so its continuous value is \(C_dr_d^{\mathrm{lin}}(P)/B\).

We first compute the distinct linear residual at a center.  For a design with signals \(x,y\), write the two measured proxies as \(W_x,W_y\).  Symmetry and conditional independence given \(U\) imply
\[
 \mathbb E(U)=\mathbb E(W_x)=\mathbb E(W_y)=0,
 \quad
 \mathbb E(UW_x)=x,
 \quad
 \mathbb E(UW_y)=y,
 \quad
 \mathbb E(W_xW_y)=xy.                                   \tag{10}
\]
The intercept coefficient is therefore zero, and the remaining normal equations are
\[
 \begin{pmatrix}1&xy\\xy&1\end{pmatrix}
 \begin{pmatrix}\beta_1\\\beta_2\end{pmatrix}
 =\begin{pmatrix}x\\y\end{pmatrix}.                     \tag{11}
\]
Since \(0<x,y<1\), the determinant \(1-x^2y^2\) is positive.  Inverting (11) gives
\[
 \beta_1=\frac{x(1-y^2)}{1-x^2y^2},
 \qquad
 \beta_2=\frac{y(1-x^2)}{1-x^2y^2}.                       \tag{12}
\]
The projection identity now yields
\[
 \begin{aligned}
 r_d^{\mathrm{lin}}(P_{\chi,\epsilon})
 &=1-(x,y)
 \begin{pmatrix}1&xy\\xy&1\end{pmatrix}^{-1}
 \begin{pmatrix}x\\y\end{pmatrix}\\
 &=\frac{(1-x^2)(1-y^2)}{1-x^2y^2}.                       \tag{13}
 \end{aligned}
\]
At the symmetric center the conditional mean \(\mathbb E(U\mid W_x,W_y)\) has the same projection error; equivalently direct enumeration of the four proxy cells gives the right side of (13).  Hence (13) equals \(R_d(P_{\chi,\epsilon})\) at these centers.  No off-center equality is used.

For \(d_{\mathrm p}\), substitute \(x=\epsilon\) and \(y^2=1-3/(10\chi)\) in (13); for \(d_{\mathrm e}\), substitute \(x=y=1/2\).  This gives
\[
 r_{d_{\mathrm p}}^{\mathrm{lin}}
 =\frac{3}{10\chi}
 \frac{1-\epsilon^2}{1-\epsilon^2(1-3/(10\chi))},
 \qquad
 r_{d_{\mathrm e}}^{\mathrm{lin}}=\frac35.               \tag{14}
\]
The first fraction multiplying \(3/(10\chi)\) is less than one because its denominator exceeds its numerator by \(3\epsilon^2/(10\chi)\).  Since \(C_{d_{\mathrm p}}=5\chi\) and \(C_{d_{\mathrm e}}=5\), (14) proves
\[
 B\Psi_{\mathrm{CLW},d_{\mathrm p}}
 =\frac32\frac{1-\epsilon^2}{1-\epsilon^2y_\chi^2}
 <\frac32<3
 =B\Psi_{\mathrm{CLW},d_{\mathrm e}}.                    \tag{15}
\]
It also proves the unrestricted-risk ordering at the center.  Thus both continuous selectors choose \(d_{\mathrm p}\).

For the exact integer rule, set \(n_d(B)=\lfloor B/C_d\rfloor\).  If \(B\ge2C_d\), the elementary inequality
\[
 0\le \frac{B}{\lfloor B/C_d\rfloor}-C_d
 \le\frac{2C_d^2}{B}                                      \tag{16}
\]
follows from \(0\le1/\lfloor z\rfloor-1/z\le2/z^2\) for \(z\ge2\).  Multiplication by the finite residual \(r_d^{\mathrm{lin}}\) shows that
\[
 B\Psi_{\mathrm{CLW},d}^{\mathrm{int}}(P;B)
 \longrightarrow C_dr_d^{\mathrm{lin}}(P).                \tag{17}
\]
The catalogue is finite and the gap in (15) is strict, so (17) gives the asserted eventual integer selector.  Equation (16) also states the positive-count boundary.

We next verify regularity.  All latent cells at the center are strictly positive.  Consistency holds by construction; experimental treatment is independent of \((U,Z(a))\); observational treatment is latent-exchangeable and has probabilities between \(3/8\) and \(5/8\); and the proxy kernels are invariant across sources.  Conditional proxy independence gives the sequential-outcome restriction.  A binary proxy with nonzero signal has a two-by-two conditional-expectation matrix with nonzero determinant.  Since \(\epsilon,y_\chi,1/2\) are nonzero, both completeness requirements and both designwise bridge operators are nonsingular.  The displayed finite bridges solve their moment equations, and all score variances are finite.  The explicit calculations below show they are strictly positive.  Thus \(P_{\chi,\epsilon}\in\mathfrak P_{\mathcal D}\); (11) is nonsingular, so it also belongs to \(\mathfrak P_{\mathcal D}^{\mathrm{CLW}}\).  Since \(Y(1)=Y(0)=U\), \(\tau=0\).  The bridge equation and source invariance give \(\mu_{d,a}=\mathbb E(U)=0\), so each bridge contrast equals \(\tau\).

For the source scores, \(h(W_y,a)=W_y/y\), so balanced experimental treatment gives
\[
 V_{E,d}=\frac4{y^2}.                                      \tag{18}
\]
The selection bridge is
\[
 q(W_x,a)=\frac{16}{15}\left\{1-\frac{(2a-1)W_x}{4x}\right\}.
\]
Expanding the four values of \((U,W_x,W_y,T)\) under the observational kernel gives
\[
 V_{O,d}=\frac{64(1-y^2)(1+14x^2)}{225x^2y^2}.             \tag{19}
\]
Substitution in (18)--(19) yields
\[
 V_{E,d_{\mathrm p}}=\frac{40\chi}{10\chi-3},
 \qquad
 V_{O,d_{\mathrm p}}
 =\frac{64}{75(10\chi-3)}(\epsilon^{-2}+14),              \tag{20}
\]
while \(V_{E,d_{\mathrm e}}=16\) and \(V_{O,d_{\mathrm e}}=384/25\).  The budget-profile theorem therefore gives
\[
 \mathcal V_{d_{\mathrm e}}^*
 =\left(\sqrt{2\cdot16}+\sqrt{3\cdot384/25}\right)^2
 =\frac{3872}{25}.                                         \tag{21}
\]
Substitution of (20)--(21) and the costs into the profiled formula gives exactly (6) of the proposed statement.  At the centers \(p_{E,1}=p_{O,1}=1/2\), hence the definition of the IKM normalization gives \(V_{g,d}^{\mathrm{IKM}}=V_{g,d}/4\).  Every profiled IKM value is one quarter of its arm-normalized counterpart, so the factor cancels in the ratio.  Since the CLW and unrestricted selectors both choose \(d_{\mathrm p}\), all four ratios agree.

Dropping the first nonnegative square-root term in the exact ratio gives
\[
 \rho\ge
 \frac{5\chi+1}{363(10\chi-3)}(\epsilon^{-2}+14)
 >\frac{\epsilon^{-2}+14}{726},                            \tag{22}
\]
because
\[
 \frac{5\chi+1}{363(10\chi-3)}-\frac1{726}
 =\frac5{726(10\chi-3)}>0.                                \tag{23}
\]
Thus \(d_{\mathrm e}\) is the unique value minimizer and the displayed lower bounds follow.

It remains to optimize over \(\chi\).  Put
\[
 A(\chi)=10\chi+1+\frac3{10\chi-3},
 \qquad
 D(\chi)=\frac{16}{75}+\frac{16}{15(10\chi-3)},
 \qquad C_0=\frac{25}{3872}.
\]
Algebra applied to the exact ratio gives
\[
 \epsilon^2\rho
 =C_0\left[\epsilon^2A(\chi)+D(\chi)(1+14\epsilon^2)
 +2\epsilon\sqrt{1+14\epsilon^2}\sqrt{A(\chi)D(\chi)}\right]. \tag{24}
\]
For \(\chi=k\epsilon^{-2/3}\), locally uniformly for \(k\) in compact subsets of \((0,\infty)\),
\[
 \epsilon^2\rho
 =\frac1{726}+C_0\epsilon^{2/3}
 \left\{2\sqrt{\frac{32k}{15}}+\frac8{75k}\right\}
 +o(\epsilon^{2/3}).                                      \tag{25}
\]
The rescaled minimizers stay in a compact subset of \((0,\infty)\): the excess in braces diverges at both endpoints, while any fixed interior \(k\) supplies a finite comparison.  Differentiation gives
\[
 q'(k)=\sqrt{\frac{32}{15}}k^{-1/2}-\frac8{75}k^{-2},
\]
which changes sign exactly once, at \(k_\star=\sqrt[3]{18}/15\).  Finally
\[
 C_0q(k_\star)=\frac{5\sqrt[3]{12}}{968}.
\]
This proves attainment for small \(\epsilon\), the optimizer scale, the sharp expansion, and every assertion.